\chapter{МГД-рівновага}\label{ch:equilibrium}

Рівновага задає геометрію магнітних поверхонь, уздовж яких рухаються потоки плазми, розглянуті в розд.~\ref{ch:transport}--\ref{ch:flows}. Нижче виведено рівняння Ґреда--Шафранова (Grad--Shafranov equation, GS), поставлено задачу з вільною межею, розібрано параметризації струму Maxfea~\cite{Merlo2011} і коду проєкту~\cite{ProjGSsolver} та реконструкцію рівноваги.

\section{Ідеальна МГД і умова рівноваги}\label{sec:03-mhd}

Однорідинна резистивна МГД у формі, яку використовує~\cite[с.~15]{Merlo2011}, складається з рівнянь Максвелла без струму зміщення, закону Ома і рівняння руху ($\rho_m$~--- густина маси):
\begin{equation}\label{eq:03-mhd}
\grad\times\vect{E}=-\pa_t\Bvec,\quad \grad\times\Bvec=\mu_0\vect{J},\quad \grad\cdot\Bvec=0,\quad
\vect{J}=\sigma(\vect{E}+\vect{v}\times\Bvec),\quad \rho_m\frac{\dd\vect{v}}{\dd t}=\vect{J}\times\Bvec-\grad p .
\end{equation}
Ідеальна МГД відповідає $\sigma\to\infty$ (тоді $\vect{E}+\vect{v}\times\Bvec=0$). У статичній рівновазі ($\vect{v}=0$, $\pa_t=0$) лишається баланс сил~\cite[с.~16]{Merlo2011}
\begin{equation}\label{eq:03-force}
\vect{J}\times\Bvec=\grad p,\qquad \grad\times\Bvec=\mu_0\vect{J},\qquad \grad\cdot\Bvec=0 .
\end{equation}
Скалярне множення \eqref{eq:03-force} на $\Bvec$ і на $\vect{J}$ дає
\begin{equation}\label{eq:03-surfaces}
\Bvec\cdot\grad p=0,\qquad \vect{J}\cdot\grad p=0 .
\end{equation}
Отже, і силові лінії, і лінії струму лежать на поверхнях $p=\text{const}$. У токамаку це вкладені тори (магнітні поверхні). Така рівновага --- вихідний стан аналізу стійкості: NOVA-W~\cite[с.~7]{Ward1992} розв'язує лінеаризовані рівняння ідеальної МГД саме навколо \eqref{eq:03-force} (розд.~\ref{ch:stability}).

\section{Рівняння Ґреда--Шафранова}\label{sec:03-gs}

В осесиметричній системі ($\pa_\tor\equiv0$) у координатах $(R,\tor,Z)$ бездивергентне поле можна записати через дві скалярні функції:
\begin{equation}\label{eq:03-B}
\Bvec=\grad\psi\times\grad\tor+F\grad\tor,\qquad \BR=-\frac{1}{R}\frac{\pa\psi}{\pa Z},\quad \BZ=\frac{1}{R}\frac{\pa\psi}{\pa R},\quad \Btor=\frac{F}{R},\quad \Bpol=\frac{|\grad\psi|}{R}.
\end{equation}
Тут $\psi$ --- полоїдальний потік \emph{на радіан} (конвенція EFIT): повний полоїдальний потік, що протікає крізь круг, обмежений колом радіуса $R$ на висоті $Z$ (за калібрування $\psi=0$ при $R=0$), дорівнює $2\pi\psi$. Функція $F=R\Btor$ пов'язана з повним полоїдальним струмом $I_{\mathrm{pol}}$ у крузі, який охоплює точка $(R,Z)$: $F=\mu_0I_{\mathrm{pol}}/2\pi$~\cite[с.~29]{Merlo2011}.

\paragraph{Узгодження конвенцій.} (i)~Merlo~\cite[с.~29--30]{Merlo2011}: $\psi$ --- повний потік, а в GS підставлено $\tilde\psi=\psi/2\pi$, тобто наш $\psi$; $f\equiv F$. (ii)~Ward і Jardin~\cite[с.~7]{Ward1992}: $\Bvec=\grad\tor\times\grad\psi+g\grad\tor$, де $2\pi\psi$ --- потік. Це наша формула із заміною $\psi\to-\psi$ і $g\equiv F$. (iii)~Нормований потік $\psiN=(\psi-\psiax)/(\psib-\psiax)$ однаковий в усіх конвенціях: множник $2\pi$ і знак скорочуються. Нормований потік Merlo $\bar\psi=(\psi_{\mathrm{ax}}-\psi)/(\psi_{\mathrm{ax}}-\psi_{\mathrm{b}})$ збігається з $\psiN$. (iv)~У коді~\cite{ProjGSsolver} (\texttt{compute.py}) $\psi$ --- «poloidal magnetic flux / 2pi», тобто теж на радіан.

\begin{derivation}[від балансу сил до рівняння GS]
\emph{Струм.} Застосуємо $\grad\times$ до \eqref{eq:03-B}. Використаємо тотожності $\grad\times(F\grad\tor)=\grad F\times\grad\tor$ і $\grad\times(\grad\psi\times\grad\tor)=-\GSop\psi\,\grad\tor$, де
\begin{equation}\label{eq:03-GSop}
\GSop\psi\equiv R\frac{\pa}{\pa R}\Big(\frac1R\frac{\pa\psi}{\pa R}\Big)+\frac{\pa^2\psi}{\pa Z^2}.
\end{equation}
Друга тотожність випливає з $\grad\tor=\hat{\vect{e}}_\tor/R$ і $\grad\cdot(R^{-2}\grad\psi)=R^{-2}\GSop\psi$. Отже,
\[
\mu_0\vect{J}=\grad F\times\grad\tor-\GSop\psi\,\grad\tor,\qquad \mu_0\Jtor=-\GSop\psi/R .
\]
\emph{Сила.} Розкладемо $\mu_0\vect{J}\times\Bvec$ на чотири доданки, використовуючи $|\grad\tor|^2=R^{-2}$ та $\grad\tor\cdot\grad\psi=\grad\tor\cdot\grad F=0$:
\[
(\grad F\times\grad\tor)\times F\grad\tor=-\frac{F\grad F}{R^2},\qquad
(\grad F\times\grad\tor)\times(\grad\psi\times\grad\tor)=-\big[(\grad F\times\grad\tor)\cdot\grad\psi\big]\grad\tor,
\]
\[
-\GSop\psi\,\grad\tor\times(\grad\psi\times\grad\tor)=-\frac{\GSop\psi}{R^2}\grad\psi,\qquad \grad\tor\times\grad\tor=0 .
\]
Права частина $\mu_0\grad p$ \eqref{eq:03-force} не має тороїдальної компоненти. Тому $(\grad F\times\grad\tor)\cdot\grad\psi=0$, тобто $\grad F\parallel\grad\psi$ і $F=F(\psi)$. З $\Bvec\cdot\grad p=0$ \eqref{eq:03-surfaces} так само $p=p(\psi)$. Тоді $\grad F=F'\grad\psi$ і $\grad p=p'\grad\psi$, а полоїдальна частина балансу набуває вигляду $-(\GSop\psi+FF')\grad\psi/R^2=\mu_0p'\grad\psi$.
\end{derivation}

Результат --- рівняння Ґреда--Шафранова~\cite[с.~30, (3.1)]{Merlo2011}:
\begin{equation}\label{eq:03-GS}
\boxed{\;\GSop\psi=-\mu_0R^2\frac{\dd p}{\dd\psi}-F\frac{\dd F}{\dd\psi}=-\mu_0R\Jtor,\qquad
\Jtor=Rp'(\psi)+\frac{FF'(\psi)}{\mu_0R}\;}
\end{equation}
Це нелінійне еліптичне рівняння. Функції $p(\psi)$ і $F(\psi)$ з МГД не визначаються: їх задають (прямий розрахунок) або підганяють під виміри (реконструкція, п.~\ref{sec:03-recon}). Тиск створює внесок у $\Jtor$, пропорційний $R$, а полоїдальний струм --- пропорційний $1/R$. На цій різниці тримаються параметризація (п.~\ref{sec:03-param}) і виродження (п.~\ref{sec:03-recon}).

\section{Задача з вільною межею}\label{sec:03-free}

Межа плазми --- це остання замкнена магнітна поверхня (LCFS). Вона або торкається лімітера, або проходить крізь сідлову X-точку (дивертор, сепаратриса)~\cite[с.~30]{Merlo2011}. Положення межі заздалегідь невідоме, тому в задачі з вільною межею \eqref{eq:03-GS} розв'язують в усій камері. Струм $\Jtor(\psi)$ відмінний від нуля в області $\psiN<1$, а поза нею є лише струми котушок $I_c$.

\begin{outofcorpus}[функція Гріна осесиметричного витка]
Розв'язок рівняння $\GSop\psi=-\mu_0R\Jtor$ у вільному просторі дається інтегралом
$\psi(R,Z)=\iint G(R,Z;R',Z')\Jtor\,\dd R'\dd Z'+\sum_cG(R,Z;R_c,Z_c)I_c$, де
\begin{equation}\label{eq:03-green}
G=\frac{\mu_0}{2\pi}\frac{\sqrt{RR'}}{k}\Big[(2-k^2)K(k^2)-2E(k^2)\Big],\qquad k^2=\frac{4RR'}{(R+R')^2+(Z-Z')^2},
\end{equation}
а $K(m)$, $E(m)$ --- повні еліптичні інтеграли з параметром $m=k^2$ (конвенція \texttt{scipy.special.ellipk}).
\end{outofcorpus}
Такі функції Гріна використовують і XLOC~\cite[с.~48]{Merlo2011}, і код проєкту~\cite{ProjGSsolver} (\texttt{GreenFunction}; до 29.09.2026 код плутав $k$ і $m=k^2$ --- E33, п.~\ref{sec:03-practice}).

\paragraph{Ітерація Пікара.} Коди з вільною межею ітерують за такою схемою. Нехай $\psi_k$ --- поточне наближення: (1)~знайти вісь (O-точку) і межу $\psib$ за $\psi_k$; (2)~обчислити $\Jtor(\psi_k)$ і нормувальні множники зі зв'язків (наприклад, $\Ip$, $\betap$); (3)~визначити $\psi$ на межі розрахункової області через функції Гріна; (4)~розв'язати лінійну задачу $\GSop\psi_{k+1}=-\mu_0R\Jtor$ і повторювати до збіжності.
Такий цикл описано в README~\cite{ProjGSsolver}. Maxfea (FORTRAN77, скінченні елементи) додатково підлаштовує струми частини PF-котушок, щоб вісь опинилася в заданій точці~\cite[с.~31--32]{Merlo2011}. SPIDER у тренажері Fenix розв'язує GS з вільною межею разом з рівняннями кіл котушок~\cite[с.~4]{Fable2022}. TSC розв'язує рівновагу на двовимірній декартовій сітці, зшиваючи плазму, вакуум і провідники за неперервністю полоїдального потоку~\cite{Jardin2000}.

\paragraph{Математична коректність.} Zou, Li і Lv~\cite[с.~2--4]{Zou2013} розглядають модельну задачу в перерізі $\Omega$ (з плоским лапласіаном замість $\GSop$):
\[
-\Delta u=f(x,u)\ \text{в}\ \{u>0\},\qquad -\Delta u=0\ \text{в}\ \{u\le0\},\qquad u|_{\pa\Omega}=\gamma<0,\qquad -\oint_{\pa\Omega}\frac{\pa u}{\pa n}=I .
\]
Плазма займає $\{u>0\}$, стала $\gamma$ невідома, повний струм $I$ задано. Закон тиску нелокальний (модель Ґреда для адіабатичної плазми): $f=\lambda_Z\,h(u_+)\,|k(|u>u(x)|)\,u'_*(|u>u(x)|)|^{q_{\mathrm Z}}+g(x)$, $1\le q_{\mathrm Z}<2$, $0\le g\in L^{r}(\Omega)$, $r>1$. Через спадну перестановку $u_*$ він нелокальний, розривний і немонотонний. Теорема~1.1 гарантує існування розв'язку $u\in H^1\cap W^{2,s_{\mathrm Z}}$, $s_{\mathrm Z}=\min\{2/q_{\mathrm Z},r\}$ (у~\cite{Zou2013} показники позначено $q$ і $\alpha$), за умов $\int_\Omega g\,\dd x>I$, монотонності $h$ (у~\cite{Zou2013} --- $\varphi$; $\lambda_Z$ --- відношення тиску частинок до магнітного) та інтегральної умови (1.3). Умова $\int g>I$ забезпечує вільну межу: $u|_{\pa\Omega}<0$, тобто плазма не торкається стінки (зауваження~1.1); теореми~1.2--1.4~\cite[с.~5]{Zou2013} послаблюють припущення. \textbf{Доведено лише існування, а не єдиність}: до цього ми повернемося в п.~\ref{sec:03-practice}.

\section{Параметризація струму і величина \texorpdfstring{$\Lambda$}{Λ}}\label{sec:03-param}

У Maxfea профіль струму задають так (параметри $\alpha,\beta$ оригіналу перейменовано на $\aJ,\bJ$):
\begin{equation}\label{eq:03-merloJ}
\Jtor(\psiN)=\lambda\Big[\bJ\frac{R}{\Rax}+(1-\bJ)\frac{\Rax}{R}\Big]\big(1-\psiN^{\aJ}\big)
\end{equation}
\cite[с.~31]{Merlo2011}. Множник $\lambda$ нормує струм на $\Ip$. Порівняння з \eqref{eq:03-GS} дає $p'=\lambda\bJ(1-\psiN^{\aJ})/\Rax$ і $FF'=\mu_0\lambda(1-\bJ)\Rax(1-\psiN^{\aJ})$. Отже, $\bJ$ --- частка струму, яку несе тиск. Відповідний тиск $p(\psiN)=p_{\mathrm{ax}}\big[1-\tfrac{\aJ+1}{\aJ}\psiN+\tfrac{1}{\aJ}\psiN^{\aJ+1}\big]$ обертається в нуль на межі~\cite[с.~32]{Merlo2011}.

Код проєкту~\cite{ProjGSsolver} має дві форми. Чисельний розв'язувач (\texttt{src/numerical}) використовує\linebreak $\lambda\big[\beta_0R/R_0+(1-\beta_0)R_0/R\big](1-\psiN^{m})^{n}$, де $R_0$ --- центр розрахункової області, а не вісь. PINN (функція \texttt{compute\_Jphi} у \texttt{GradShafranov.py}) розділяє показники для двох доданків:
\begin{equation}\label{eq:03-pinnJ}
\Jtor=\lambda\Big[\beta_0\,\hat R\,\big(1-\psiN^{\alpha_m}\big)^{\alpha_n}+(1-\beta_0)\,\hat R^{-1}\big(1-\psiN^{\beta_m}\big)^{\beta_n}\Big],\qquad \hat R=R/R_c,
\end{equation}
із типовими значеннями $\alpha_m=\beta_m=2$, $\alpha_n=\beta_n=1$, $\beta_0=0{,}5$, $R_c=\SI{1.8}{м}$ (\texttt{train\_PINN.py}; $\alpha_{m,n},\beta_{m,n},\beta_0$ --- імена змінних коду, а не $\beta$ плазми). $\lambda$ і $\beta_0$ оновлюються зі зв'язків на $\Ip$ і $\betap$ (\texttt{profiles.py}).

\paragraph{Вертикальне поле і $\Lambda$.} Параметри $\aJ,\bJ$ добирають так, щоб рівновага відтворювала виміряні магнітні сигнали. Для кругової плазми радіуса $r_a$ радіальне розширення компенсує вертикальне поле~\cite[с.~35]{Merlo2011}
\begin{equation}\label{eq:03-Bv}
B_{v,\mathrm{eq}}=-\frac{\mu_0\Ip}{4\pi R_0}\Big(\ln\frac{8R_0}{r_a}+\Lambda-\frac12\Big),\qquad \Lambda=\betap+\frac{\li}{2}-1,
\end{equation}
де $\betap=\fsa{p}/(B_{\mathrm{p}}^2(r_a)/2\mu_0)$ (у Merlo $\beta_\theta$), а $\li$ --- нормована внутрішня індуктивність~\cite[с.~33]{Merlo2011} ($\Lambda$ не плутати з $\lnL$). Величину $\Lambda$ можна оцінити з вимірів: з $B_{\mathrm{p}}$ на 8 зондах при $r=r_b$ беруть середнє $B_{\mathrm{p},0}$ і косинусну гармоніку $B_{\mathrm{p},c}$. Тоді $\Lambda^*=(B_{\mathrm{p},c}/B_{\mathrm{p},0})(R_0/r_b)$, і з тороїдальною поправкою
$\Lambda=2(1+r_a^2/r_b^2)^{-1}\big[\Lambda^*+\tfrac14(1-r_a^2/r_b^2)-\tfrac12\ln(r_b/r_a)\big]$~\cite[с.~35]{Merlo2011}.

\begin{corpusexample}[RFX-mod у режимі токамака~\cite{Merlo2011}]
Параметри установки: $R_0=\SI{2}{м}$, радіус графітової першої стінки \SI{0.459}{м}, $\Ip\le\SI{200}{кА}$ (с.~19--20).
Для кругової плазми (розряд \#29648) спершу взято $\aJ=3$, $\bJ=0{,}25$, що дало $\li=0{,}76$, $\betap=0{,}26$ (с.~33). Для кращої узгодженості з виміряним $\Lambda$ перейшли до $\aJ=1$, $\bJ=0{,}25$ (с.~38): $\li=1{,}00$, $\betap=0{,}26$ для кругової плазми і $\li=0{,}95$, $\betap=0{,}29$ для double-null (\#29746).
Резистивну плазму моделюють спадом потоку на осі $2\pi\,\dd\psi_{\mathrm{ref}}/\dd t=-R_{\mathrm{pla}}\Ip$ (у~\cite{Merlo2011} $\psi$ --- повний потік) з $R_{\mathrm{pla}}=\SI{12.9}{\micro\ohm}$ (с.~33--34).
\end{corpusexample}

\section{Зсув Шафранова і реконструкція рівноваги}\label{sec:03-recon}

\begin{outofcorpus}[зсув Шафранова]
Тиск і тороїдальність зсувають магнітні поверхні назовні. Для кругових поверхонь з великим аспектним відношенням зсув центру поверхні радіуса $r$ задовольняє рівняння
\[
\frac{\dd\Delta_{\mathrm{S}}}{\dd r}=-\frac{r}{R_0}\Big[\betap(r)+\frac{\li(r)}{2}\Big],\qquad \Delta_{\mathrm{S}}(a)=0,
\]
де $\betap(r)$ і $\li(r)$ обчислено для поверхні радіуса $r$. Звідси $\Delta_{\mathrm{S}}(0)\sim(a^2/2R_0)(\betap+\li/2)$. Це та сама комбінація $\betap+\li/2=\Lambda+1$, що й у \eqref{eq:03-Bv}.
\end{outofcorpus}

\paragraph{Експериментальна геометрія поверхонь.} На TFTR геометрію поверхонь отримали з багатоканальної FIR-інтерферометрії методом \emph{асиметричної інверсії Абеля}. Метод~\cite[с.~3]{Park1989} відмовляється від припущення про симетрію профілю і нульову густину на краях, а зсунуті поверхні визначає ітераціями разом із магнітними вимірами.

\paragraph{Магнітна реконструкція.} Коди реального часу розкладають $\psi$ за базисом~\cite[с.~48--52]{Merlo2011}. XLOC на JET використовує поліноми від $\xi=R^2-R_0^2$ і $Z$ (до сьомого степеня). Базис накладає вакуумну умову $\GSop\psi=0$, коефіцієнти знаходять методом найменших квадратів. Переріз ділять на зони зі «зшиванням», а внесок диверторних котушок віднімають і додають назад через функції Гріна. Спрощений метод Merlo інтерполює 24 значення потоку (8 виміряних і 16 екстрапольованих) тонкопластинним сплайном з базисною функцією $r_{\mathrm s}^2\ln r_{\mathrm s}^2$ ($r_{\mathrm s}$ --- відстань до вузла). Рівняння GS при цьому не накладається; параметр згладжування $\nu=0$ для даних симулятора і $\nu\approx0{,}04$ для експерименту.
\begin{outofcorpus}[EFIT]
EFIT (Lao та ін., 1985) подає $p'(\psiN)$ і $FF'(\psiN)$ як лінійні комбінації базисних функцій. За фіксованого $\psi$ підгонка коефіцієнтів до вимірів є лінійною задачею найменших квадратів, і цикл чергує її з розв'язком \eqref{eq:03-GS}.
\end{outofcorpus}

\paragraph{Виродження $p'$ проти $FF'$.} З \eqref{eq:03-GS} видно, що $p'$ і $FF'$ розрізняються лише через зміну $R^2$ уздовж магнітної поверхні, тобто на рівні $\order{a/R_0}$. Для кругової плазми зовнішні магнітні виміри дають лише комбінацію $\Lambda$, а не окремо $\betap$ і $\li$~\cite[с.~35]{Merlo2011}. Fenix відтворив струми PF-котушок AUG лише тоді, коли крайовий бутстреп-струм зменшили так, щоб узгодити $l_{i3}$ з реконструкцією~\cite[с.~6--7]{Fable2022}. Отже, струми котушок чутливі до профілю струму; наскільки сильно це виродження псує реконструкцію без внутрішніх діагностик (MSE, поляриметрія), --- відкрите питання \textbf{Q2}~\cite{ProjRegisters}.

\section{Практикум і дослідницький трек}\label{sec:03-practice}


\begin{practicum}[розв'язувач GS з вільною межею: відтворити $\to$ перевірити тести $\to$ розширити]
\textbf{Крок 0 (відтворення).} За README~\cite{ProjGSsolver}: \texttt{python3 -m venv .venv}, \texttt{pip install -r requirements.txt} (або \texttt{conda env create -f environment.yaml}), далі\\
\texttt{python3 test\_free\_boundary.py}\\
Типовий приклад (сітка $128\times128$, $\Ip=\SI{200}{кА}$, дві котушки по $-\SI{90}{кА}$ у $(2{,}2;\,\pm0{,}7)$~м) збігається за 116 ітерацій Пікара приблизно за \SI{0.55}{с} і дає рис.~\ref{fig:03-gsfb}. До 29.09.2026 ця команда падала (модуль переїхав у \texttt{src/numerical/}, \texttt{SyntaxError} у \texttt{grad\_shafranov.py}, невизначений \texttt{GSMatrix}, не було \texttt{environment.yaml}), а сам розв'язувач з вільною межею не міг дати фізичного результату: $\psi$ на межі бралося як одне число з двох країв, маски плазми не було, осі $R$ і $Z$ на рисунку були переставлені. Ремонт (E34~\cite{ProjRegisters}) задокументовано у файлі \texttt{CHANGES-2026-09-29.md}.

\textbf{Крок 1 (тести).} \texttt{python3 -m pytest} --- 22 тести. \texttt{tests/test\_green.py} порівнює $G$ \eqref{eq:03-green} з прямим інтегруванням $\psi=RA_\tor$ по витку (Біо--Савар, квадратура; відхилення $\sim10^{-15}$) і перевіряє симетрію $G(R,Z;R',Z')=G(R',Z';R,Z)$. Такий тест і виявив помилку E33: код передавав у \texttt{ellipk}/\texttt{ellipe} модуль $k$ замість параметра $m=k^2$. Похибка старої $G$: $+9\,\%$ біля нитки, $+42\,\%$ при $k^2=0{,}95$ (виток \SI{1}{А} у $(1{,}5;\,0{,}25)$~м, точка $(1{,}0;\,0)$~м: $3{,}355\cdot10^{-7}$ замість $2{,}355\cdot10^{-7}$~Вб/рад), ${\times}8$ при $k^2=0{,}5$; при $k\to0$ вона розбігається, тож потік далеких котушок був хибним на порядки~\cite{ProjRegisters}. Та сама помилка сиділа в PINN-копії, тому всі результати PINN, отримані до виправлення, використовували хибну $G$. \texttt{tests/test\_free\_boundary\_smoke.py} перевіряє, зокрема, повний струм трьома способами, серед них законом Ампера $-\oint\Bvec_{\mathrm p}\cdot\dd\vect l/\mu_0$ довкола області, який не залежить від дискретизації. \emph{Завдання:} поверніть у $G$ модуль $k$ замість $k^2$ і переконайтеся, що тести ловлять помилку.

\textbf{Крок 2 (розширення).} Обмеження, що лишились, --- це завдання. (а)~LCFS обмежує прямокутник сітки (рис.~\ref{fig:03-gsfb}): додайте пошук X-точки (сідлова точка $\grad\psi=0$) і контур лімітера. (б)~Без зворотного зв'язку за положенням ітерація Пікара нестійка для багатьох наборів котушок (плазма дрейфує до стінки); додайте підлаштування струмів котушок, як у Maxfea. (в)~Додайте зв'язок на $\betap$. (г)~Шлях із фіксованою межею (\texttt{test\_boundary.py}) досі падає (\texttt{ValueError}: магнітну вісь не знайдено) --- полагодьте його. (д)~PINN без даних KSTAR (їх немає в репозиторії) не навчено й не перевірено.

\textbf{Що виміряти:} $\|\psi_{k+1}-\psi_k\|/\|\psi_k\|$ залежно від $k$; положення осі й LCFS залежно від струмів котушок; нев'язку $\GSop\psi+\mu_0R\Jtor$ (має спадати як $h^2$ зі зменшенням кроку сітки $h$); GS-неузгодженість з блоку E6--E8 для отриманого $\psi$.
\end{practicum}

\begin{figure}[htbp]\centering
\includegraphics[width=0.9\textwidth]{gs_free_boundary_fixed.jpg}
\caption{Рівновага з вільною межею після ремонту \texttt{Tokamak-GS-solver} (29.09.2026), типовий приклад \texttt{test\_free\_boundary.py}: сітка $128\times128$ у прямокутнику $R\in[1;\,2]$~м, $Z\in[-0{,}5;\,0{,}5]$~м; $\Ip=\SI{200}{кА}$ (множник $\lambda$ підлаштовується під $\Ip$), $\beta_0=1$, $m=n=1$; дві котушки PF по $-\SI{90}{кА}$ у $(2{,}2;\,\pm0{,}7)$~м поза сіткою. Ліворуч --- $\psi$ [Вб/рад], хрестик --- магнітна вісь ($R\approx\SI{1.42}{м}$), червоний контур --- LCFS ($\psiN=1$); праворуч --- $\Jtor$ [А/м$^2$], нуль поза маскою плазми. LCFS торкається внутрішнього краю сітки $R=\SI{1}{м}$: межу плазми тут задає прямокутник розрахункової області, а не X-точка чи лімітер. Власний розрахунок~\cite{ProjGSsolver}.}\label{fig:03-gsfb}
\end{figure}

\begin{figure}[htbp]\centering
\includegraphics[width=0.85\textwidth]{d2_branch_mean.jpg}
\caption{D2: задача Брату (на рисунку $\lambda\equiv\lambda_{\mathrm{B}}$). (a)~Дві точні гілки при $\lambda_{\mathrm{B}}=1$ і їхнє $L^2$-середнє. (b)~Нев'язка $u''+\lambda_{\mathrm{B}}\ee^u$: для гілок нуль, для середнього мінімум $\approx-22$ у центрі. (c)~Відносна нев'язка середнього і гілки та відстань між гілками залежно від $\lambda_{\mathrm{B}}$; пунктир --- точка складки; контроль Біна --- на рівні машинної точності. Власний розрахунок~\cite{ProjDeepDives}.}\label{fig:03-d2}
\end{figure}

\begin{established}{E21, E24, E4, E5}
\textbf{Джерела~\cite{ProjRegisters}.} E21: розв'язок GS \emph{може бути неєдиним} на реальній геометрії (Ham і Farrell, 2024; Pentland та ін., 2025, MAST-U, дефляційна континуація). E24: критичний стан Біна опуклий і має єдиний розв'язок, а задача плазми з вільною межею неопукла і багатозначна.
\textbf{Власний вимір~\cite{ProjDeepDives} (D2).} Як сурогат механізму взято одновимірну задачу Брату $-u''=\lambda_{\mathrm{B}}\ee^{u}$, $u(0)=u(1)=0$. Нижче точки складки $\lambda_{\mathrm{B}}^*=3{,}513830719$ вона має рівно дві гілки; обчислена $\lambda_{\mathrm{B}}^*$ збігається з літературою до останньої цифри (E5). E4: $L^2$-середнє двох гілок (до нього збігається регресор з $L^2$-втратою) має відносну нев'язку $1{,}22$ проти $6{,}4\cdot10^{-8}$ у гілки; в опуклому контролі (Бін) нев'язка середнього $2{,}5\cdot10^{-13}$.
\textbf{Застереження:} Брату --- \emph{не} GS; перенесення на GS спирається на цитати (рис.~\ref{fig:03-d2}).
\end{established}



\begin{conjecture}{C7}
\textbf{Що стверджує:} Deep Ritz (мінімізація функціонала енергії нейромережею) на неопуклій задачі плазми буде нестабільним відносно seed (обиратиме різні гілки). На опуклій задачі Біна він буде стабільним~\cite{ProjRegisters}. \textbf{Що спростує:} обидві задачі стабільні, тобто неопуклість не виявляється на доступних масштабах. \textbf{Вартість:} близько дня, потрібні обидва розв'язувачі. \textbf{Що зробити:} взяти \texttt{src/examples/deep\_ritznet.py}~\cite{ProjGSsolver} (опуклий пуассонівський функціонал) і замінити функціонал на $\int(\tfrac12u'^2-\lambda_{\mathrm{B}}\ee^u)$. Прогнати 20 seed і класифікувати результати за $\max u$. Увага: функціонал необмежений знизу, а верхня гілка --- сідло, тож «нестабільність» може виявитися розбіжністю, а не вибором іншої гілки; ці випадки треба розрізняти.
\end{conjecture}


\begin{figure}[htbp]\centering
\includegraphics[width=\textwidth]{d1_amplification.jpg}
\caption{D1: медіанна похибка семи скалярів (в одиницях популяційного $\sigma$) залежно від відносної похибки $e=\|\delta\psi\|/\|\psi-\bar\psi\|$, $R^2_\psi=1-e^2$. Ліворуч --- білий шум, праворуч --- хвіст PCA. Усі нахили $a_e<1$. 60 кадрів трьох демо-розрядів DIII-D. Власний розрахунок~\cite{ProjDeepDives}; дані:~\cite{FEC2026}, CC~BY~4.0.}\label{fig:03-d1}
\end{figure}

\begin{established}{E6--E8}
Нев'язку GS можна обчислити з одного лише $\psi$, без поданих $p'$ і $FF'$, бо в структурі EFIT джерело лінійне за коефіцієнтами профілів~\cite{ProjRegisters}:
\[
g(\psi)=\min_{c_i,d_j}\frac{\big\|\GSop\psi-\sum_{i<4}c_iR^2\psiN^{\,i}-\sum_{j<4}d_j\psiN^{\,j}\big\|}{\|\GSop\psi\|}
\]
Норму беруть по пікселях плазми (\texttt{Our try/04-novelty/gs\_residual\_probe.py}). На 12 кадрах DIII-D~\#203702: E6 --- істина (EFIT) $0{,}009$; з 1\,\% шуму $0{,}650$ (у 70 разів більше); з 3\,\% шуму $0{,}921$; E7 --- $g$ інваріантна до масштабу: $\psi\times1{,}05$ дає те саме значення; E8 --- $g$ сліпа до гладких похибок: гаусове згладжування $0{,}010$, PCA-5 $0{,}0093$ (раніше в реєстрі стояло $0{,}016$ без скрипта-джерела; перерахунок 2026-09-29 тими самими функціями дає $0{,}0093$, тож висновок про сліпоту лише посилився). Причина сліпоти: $\GSop$ --- оператор другого порядку, він підсилює високі хвильові числа.
\end{established}


\begin{established}{E1--E3, E14, E36}
Скаляри рівноваги є диференціальними функціоналами $\psi$~\cite{ProjRegisters,ProjDeepDives}: вісь --- критична точка $\grad\psi=0$, форма (витягнутість, трикутність, об'єм) --- контур крізь сідло, $\li$ --- інтеграл від $|\grad\psi|^2$. E1: відображення $\psi\to$ скаляри не ліпшицеве в $L^2$: похибка скаляра зростає як $e^{a_e}$ з $a_e\in[0{,}35;\,0{,}70]$ (у реєстрі~--- $\alpha$; рис.~\ref{fig:03-d1}). Похибка $R_{\mathrm{ax}}$ сягає $1\sigma$ вже при $R^2_\psi=0{,}99993$, а при $R^2_\psi=0{,}99$ становить 3,4~см. Це виміряно для \emph{білого} (повузлового) шуму; для гладкого збурення $a_e\approx0{,}69$--$1{,}08$ (уточнення 2026-09-30). E2: на трьох демо-розрядах показники утворюють три кластери: $0{,}53$--$0{,}55$, $0{,}35$--$0{,}42$ і $0{,}70$; на ширшій вибірці їх немає (R10). E3: хвіст PCA шкодить більше, ніж білий шум тієї самої енергії: для $Z_{\mathrm{ax}}$ $64{,}4\sigma$ проти $6{,}3\sigma$ при $e=0{,}30$. E14 (синтетичний тест): зсув осі під шумом пропорційний $1/\sqrt{\lambda_H}$, де $\lambda_H$ --- кривизна (власне значення гесіана $\psi$) в O-точці. Добуток $|\Delta x|\sqrt{\lambda_H}$ сталий з точністю $\pm7\,\%$ на 16-кратному діапазоні $\lambda_H$; тест відновлено 2026-09-30 (для $R_{\mathrm{ax}}$ $\pm5{,}6\,\%$, \texttt{c1/e14\_curvature.py}). \textbf{E36:} скорер шукає вісь як максимум на вузлах із параболічним уточненням, тому для шуму амплітуди $\eta$ у вузлі $a_e=1$, поки $\eta<\eta^{*}=\lambda_H h^2/2$, і $a_e\approx0{,}4$--$0{,}5$ вище, де дискретний максимум стрибає ($h$ --- крок сітки).
\end{established}
\begin{refuted}{R6}
«$a_e=\tfrac12$ --- універсальний закон» для положення критичної точки (осі). Спростовано: показник залежить від кривизни ($0{,}687$ при $\lambda_H=1$ і $0{,}841$ при $\lambda_H=4$)~\cite{ProjRegisters}. Причину пояснює E36: у фіксованому вікні $\eta$ більша кривизна означає більшу частку лінійного режиму.
\end{refuted}

\begin{refuted}{R10 (була гіпотеза C1)}
«Три кластери E2 відповідають трьом механізмам --- положенню екстремуму, контуру крізь сідло та інтегралу». Критерії записали до перевірки. Розбиття стабільне лише на трьох демо-розрядах (96,7\,\% бутстрепу по розрядах); на 28 розрядах DIII-D --- 0\,\% ($\li$: $a_e$ з $0{,}70$ падає до $0{,}47$), на MAST --- 3,3\,\%, на 12 однонульових рівновагах Серфона--Фрайдберга~\cite{Cerfon2010} --- 0,3\,\%~\cite{ProjRegisters,ProjGuide}. Урок: кластер, побачений на трьох розрядах, --- ще не закон.
\end{refuted}

\begin{conjecture}{C3}
Кривизна $\lambda_H$ з E14 корисна для стратифікації тестової вибірки. \textbf{Спростує:} вузький розподіл $\lambda_H$ по корпусу. \textbf{Стан:} на 28 розрядах DIII-D $\lambda_H\in[0{,}13;\,3{,}6]$ --- розподіл широкий, але $a_e$ осі між терцилями $\lambda_H$ змінюється слабко~\cite{ProjRegisters}. \textbf{Що зробити:} стратифікувати за $\lambda_H$ залишки справжнього бейзлайну PCA+Ridge (код --- \texttt{Our try/03-deep-dives/D1-psi-to-scalars/c1/real\_sweep.py}). Дані --- Fusion Equilibrium Challenge~\cite{FEC2026}.
\end{conjecture}

\paragraph{Відкриті питання~\cite{ProjRegisters}.} Крім \textbf{Q2} (п.~\ref{sec:03-recon}), \textbf{Q3}: чи трапляється неєдиність GS (E21) у робочих режимах (від цього залежить, чи D2 реальна).

\subsection*{Контрольні запитання}
\begin{enumerate}
\item Виведіть \eqref{eq:03-surfaces} з \eqref{eq:03-force}. Чому лінії поля й струму лежать на поверхнях $p=\text{const}$?
\item Як переписати формули Merlo і Ward для нашого $\psi$? Чи змінюється при цьому $\psiN$?
\item Навіщо функція Гріна в задачі з вільною межею? Що забезпечує умова $\int g>I$?
\item Чому магнітні виміри кругової плазми дають лише $\Lambda$? Як це пов'язано з виродженням?
\item Чому $L^2$-середнє гілок не є розв'язком? До чого сліпа GS-нев'язка?
\end{enumerate}

\subsection*{Задачі}
\begin{enumerate}
\item \emph{Вертикальне поле RFX-mod.} Візьміть $R_0=\SI{2}{м}$, $r_a=\SI{0.459}{м}$ (радіус першої стінки), $\Ip=\SI{150}{кА}$ (порядок величини з рис.~3.2 у~\cite{Merlo2011}; перевірте: за $\Bzero=\SI{0.6}{\tesla}$ маємо $q(a)\approx2{,}1$). (а)~Обчисліть $B_{v,\mathrm{eq}}$ \eqref{eq:03-Bv} для двох параметризацій розряду \#29648 (п.~\ref{sec:03-param}); наскільки у відсотках відрізняються результати? (б)~Оцініть $\Delta_{\mathrm{S}}(0)$ за формулою з рамки для $\betap=0{,}26$, $\li=1{,}00$.
\item \emph{Брату.} Розв'язок: $u=-2\ln\bigl[\cosh\bigl(\theta(x-\tfrac12)/2\bigr)/\cosh(\theta/4)\bigr]$, $\lambda_{\mathrm{B}}=\theta^2/\bigl[2\cosh^2(\theta/4)\bigr]$.
  (а)~Знайдіть $\lambda_{\mathrm{B}}^*$. (б)~Для $\lambda_{\mathrm{B}}=1$ знайдіть $u_{\max}$ обох гілок. (в)~Обчисліть нев'язку $L^2$-середнього в точці $x=\tfrac12$ і порівняйте з рис.~\ref{fig:03-d2}(b).
\item \emph{Інваріантність E7.} Покажіть, що коли $\psi$ розв'язує \eqref{eq:03-GS} з $(p',FF')$, то $c\psi$ розв'язує його з перемасштабованими профілями. Чому \emph{незалежний} піксельний шум це руйнує?
\end{enumerate}
